\chapter{Varicocelectomy}

\section*{Introduction \& Clinical Indications}
Varicocelectomy is the ligation or occlusion of dilated veins in the spermatic cord. It may be considered for a palpable varicocele in a man attempting conception who has infertility and abnormal semen parameters, for persistent pain attributed to the varicocele, or for testicular growth or size concerns in selected adolescents. A nonpalpable varicocele found only on imaging, normal semen parameters, or an unrelated pain syndrome generally favors observation. Repair is not automatically required before assisted reproduction.

\section*{Relevant Surgical Anatomy}
The left testicular vein usually drains to the left renal vein and the right to the inferior vena cava. The spermatic cord also contains the testicular artery, deferential and cremasteric vessels, vas deferens, lymphatics, and nerves. Preserving the artery and lymphatics reduces the risks of testicular atrophy and hydrocele. The inguinal and subinguinal approaches expose different levels of the cord; microsurgical identification can reduce missed veins and injury.

\section*{Preoperative Workup \& Preparation}
Diagnosis is primarily by history and examination, with Doppler ultrasound used when the examination is uncertain, for recurrence, or to assess another scrotal condition. For infertility, semen analyses and evaluation of both partners guide counseling. Hormonal testing is selective. Consent covers observation, embolization, assisted reproduction, uncertain fertility benefit, recurrence, hydrocele, hematoma, chronic pain, and injury to the testicular artery or vas. Active infection and uncorrected coagulopathy should be treated before elective repair; anesthesia-specific fasting and antibiotic plans are individualized.

\section*{Surgical Technique \& Procedure}
\begin{enumerate}
  \item The patient is positioned supine and receives local, regional, or general anesthesia as appropriate.
  \item Through an inguinal or subinguinal incision, the spermatic cord is exposed and the testicular artery, lymphatics, vas, and veins are identified.
  \item Dilated veins are ligated or divided while preserving the artery and lymphatics. Magnification and Doppler may assist identification.
  \item Hemostasis is confirmed, the cord is returned to position, and the incision is closed.
\end{enumerate}

\section*{Complications \& Risks}
Risks include hematoma, infection, hydrocele, recurrence or persistence, injury to the testicular artery with possible testicular atrophy, chronic scrotal or groin pain, and anesthesia complications. Improvement in pain, semen parameters, or testicular growth is not guaranteed.

\section*{Postoperative Care \& Recovery}
Provide analgesia, scrotal support, wound instructions, and activity guidance. Ambulation is encouraged as tolerated. Fever, severe or increasing pain, expanding swelling, or testicular color change requires prompt assessment. Follow-up evaluates healing, pain, testicular size, and recurrence. When fertility is the indication, semen analysis is commonly repeated after about 3 months or later, with continued discussion of assisted reproduction when appropriate.

\section*{Outcomes \& Prognosis}
Fertility and semen outcomes vary with baseline semen parameters, female-partner factors, varicocele grade, and technique. Repair may improve semen parameters and pregnancy probability in selected couples, but benefit is not universal and should not be expressed as one fixed success rate.

The fertility discussion should include the expected time course and alternatives. Sperm production takes several weeks, so an early postoperative semen test may not reflect the final result. If the semen analysis remains abnormal, the couple may discuss repeat examination, endocrine or genetic evaluation, assisted reproduction, or treatment of another male or female factor. A varicocele repair should not delay evaluation when the female partner's age or ovarian reserve makes time important.

The operation is usually performed with the patient supine and a small groin incision. A microscope or loupe magnification helps distinguish the testicular artery, lymphatics, vas, and small veins. Intraoperative Doppler can support artery identification but does not replace anatomical judgment. Venous channels are divided only after adequate arterial and lymphatic preservation is established. The surgeon should avoid excessive traction on the cord and should document whether the repair was unilateral or bilateral.

Postoperative swelling can mimic recurrence during the early healing period. A persistent or recurrent dilated vein is assessed after the edema has resolved, often with standing examination and Doppler ultrasound. Hydrocele develops when lymphatic drainage is disrupted; a tense hydrocele or progressive testicular enlargement may require later treatment. Testicular atrophy is uncommon with careful technique but is a serious complication that should be included in consent. New severe pain, fever, erythema, or rapidly increasing scrotal swelling warrants urgent assessment for infection, hematoma, torsion, or vascular injury.

\section*{Counseling and Long-Term Follow-Up}
Varicocele repair is an elective decision in most fertility and pain cases. The clinician should explain that a palpable varicocele is an association, not proof that it is the cause of abnormal semen parameters or discomfort. A detailed history includes prior scrotal surgery, trauma, infection, torsion, childhood testicular position, medications, heat exposure, tobacco, alcohol, recreational drugs, anabolic steroids, and systemic illness. Testicular examination records volume, consistency, position, and the presence of a contralateral lesion. In adolescents, growth over time is often more informative than a single size difference.

Microsurgical subinguinal repair is commonly favored because it allows identification of small veins while preserving arteries and lymphatics, but the best approach depends on the surgeon, anatomy, prior operations, and patient preference. Laparoscopic, inguinal, and radiologic embolization approaches remain options. The consent discussion should include the possibility of persistent reflux, missed collateral veins, hydrocele, hematoma, infection, testicular ischemia, chronic pain, and the need for a repeat procedure. A technically uncomplicated operation cannot guarantee semen improvement or conception.

When pain is the indication, the clinician should first exclude torsion, epididymitis, hernia, tumor, referred pain, and pelvic or neuropathic causes. The pain should be compatible with a varicocele, such as a dull ache worsened by standing or exertion, rather than acute severe pain. Even in a suitable patient, pain can persist or change character after repair. Conservative measures and a period of observation may be reasonable when symptoms are mild, while persistent disabling pain requires shared decision-making about repair and its alternatives.

Follow-up includes wound and scrotal examination, assessment of testicular size and pain, and semen testing when fertility is the goal. Pregnancy is a couple-level outcome and should not be attributed to the operation without considering ovulation, tubal factors, sperm DNA or other male factors, and the time interval. Patients should be advised to return for a new hard testicular mass, sudden pain, fever, increasing swelling, or a recurrent prominent vein. Long-term fertility care may include repeat semen analysis, reproductive endocrinology, or assisted reproduction, depending on the couple's goals.

\section*{Patient Education}
Scrotal support, ice used with skin protection, and activity modification can reduce early discomfort. Patients should avoid heavy lifting until advised and should not judge fertility benefit by early postoperative symptoms. A sudden severe testicular pain, rapidly enlarging hematoma, fever, wound drainage, or a new hard testicular lump requires prompt evaluation.

\section*{Evidence and Shared Decision-Making}
Guideline recommendations emphasize a clinically palpable varicocele, infertility, and abnormal semen parameters as the setting in which repair is most likely to be considered. The quality of evidence for pregnancy and live-birth benefit is variable, so counseling should distinguish improvements in semen measures from the outcomes that matter to the couple. A normal semen analysis, a subclinical varicocele, or pain with an atypical pattern generally shifts the balance toward observation or another evaluation.

\section*{Documentation and Safety}
Documentation includes the side treated, examination findings, testicular volume, semen results when relevant, approach, artery and lymphatic preservation, veins divided, complications, and follow-up plan. Patients should receive clear instructions about scrotal support, activity, semen testing, and emergency symptoms. A recurrent vein, hydrocele, or persistent pain should prompt reassessment rather than automatic repeat surgery.

\section*{Practical Follow-Up}
Early review checks the incision, scrotal swelling, pain, and testicular examination. Semen analysis is delayed until a new spermatogenic cycle has occurred, and pregnancy counseling includes both partners. A new mass, sudden pain, fever, or rapidly increasing swelling requires urgent assessment.

\section*{Safety Summary}
The most appropriate candidates have a palpable varicocele and a clinically relevant fertility or pain indication. Preservation of the testicular artery and lymphatics is central, and a technically successful repair cannot guarantee pregnancy, pain relief, or normalized semen parameters.

\section*{Final Safety Note}
Repair is considered for selected palpable varicoceles with a relevant fertility or pain indication. It cannot guarantee pregnancy, normalized semen parameters, or pain resolution, and persistent or recurrent symptoms need reassessment.

\section*{Completion Checklist}
Document the side, palpable grade, testicular volume, fertility or pain indication, artery and lymphatic preservation, and semen follow-up. Persistent pain, a new mass, fever, or rapidly increasing swelling requires assessment.

\section*{Handoff}
The follow-up plan should state when semen testing is due, how pain will be reviewed, and which scrotal findings require urgent assessment.

\section*{References}

\section*{Additional Clinical Considerations}
The physical examination is central to selection. A clinically palpable varicocele that becomes more prominent with standing or Valsalva is different from a small venous dilation seen only on ultrasound. A right-sided, sudden-onset, non-reducible, or isolated varicocele requires assessment for retroperitoneal or renal pathology rather than routine repair. Doppler ultrasound can document venous reflux and testicular volume, but it should answer a clinical question rather than replace examination.

In an infertility evaluation, both partners should be assessed and semen abnormalities should generally be confirmed on repeat testing because results fluctuate with illness, abstinence interval, fever, and laboratory methods. Repair is most often discussed for a palpable varicocele in a man attempting conception with abnormal semen parameters and no better explanation for infertility. It is not routinely recommended for a subclinical varicocele or a man with normal semen analysis. Assisted reproduction remains an alternative, particularly when female-partner factors, severe male-factor infertility, or time constraints make surgery less suitable.

For adolescents, serial examination and testicular-volume assessment help identify persistent asymmetry or impaired growth. A single measurement can be affected by examiner and ultrasound variation. Pain attributed to a varicocele should be typical, persistent, and assessed after other causes of scrotal pain have been considered. Pain may persist after technically successful repair, and the possibility of chronic pain or a different diagnosis should be addressed before surgery.

Microsurgical subinguinal repair aims to preserve the testicular artery and lymphatics while interrupting refluxing veins, including external spermatic or gubernacular veins when they are relevant to the chosen approach. Lymphatic preservation reduces hydrocele risk. Embolization is an alternative for selected patients, but it has its own technical failure, contrast, radiation, and recurrence considerations. The operative report should document the approach, artery preservation or Doppler confirmation, veins treated, complications, and whether a hydrocele was present.

Semen parameters are not expected to change immediately. Follow-up commonly waits at least one spermatogenic cycle, and pregnancy outcomes may take longer and depend on both partners. A persistent or recurrent varicocele, hydrocele, new testicular asymmetry, fever, or worsening pain requires reassessment. Even when semen improves, the couple should be counseled that repair does not guarantee pregnancy and does not eliminate the need for timely reproductive planning.

\section*{References}
\begin{enumerate}
  \item American Urological Association and American Society for Reproductive Medicine. Diagnosis and Treatment of Infertility in Men: AUA/ASRM Guideline.
  \item European Association of Urology. Guidelines on Sexual and Reproductive Health: Male Infertility and Varicocele.
  \item World Health Organization. WHO Laboratory Manual for the Examination and Processing of Human Semen. 6th ed.
\end{enumerate}

\clearpage
